%! Linear Regression — Unit 1, Lesson 1.5 — Homework (Answer Key)
\documentclass[10pt]{article}
\usepackage{algebra2-article}
\usepackage{algebra2-key}   % pulls in -boxes; do NOT also load -boxes
\newcommand{\TallMath}[1]{$\displaystyle #1\rule[-1.4em]{0pt}{3.2em}$}
% Coordinate grid: {xmin}{xmax}{ymin}{ymax}
\newcommand{\lgrid}[4]{%
  \draw[step=1, linegray!40, very thin] (#1,#3) grid (#2,#4);
  \draw[->, semithick, charcoal] (#1,0)--(#2,0) node[below right, font=\scriptsize]{$x$};
  \draw[->, semithick, charcoal] (0,#3)--(0,#4) node[left, font=\scriptsize]{$y$};}
\newcommand{\dpt}[2]{\fill[navy] (#1,#2) circle (2.8pt);}

\begin{document}

\pageheader{Unit 1, Lesson 1.5}{Homework --- Answer Key}

\vspace{0.10in}

\begin{notesbox}{Practice}
Show your work. Describe an association by \emph{direction}, \emph{form}, and \emph{strength}. The
correlation coefficient $r\in[-1,1]$ gives direction (sign) and strength ($|r|$). To predict with a line of
best fit $\hat y=ax+b$, substitute an $x$.

\begin{enumerate}[leftmargin=1.9em, itemsep=11pt]
  \item \textbf{Classify the direction} (positive / negative / none) of each scatterplot.
        \begin{center}
        \begin{tikzpicture}[scale=0.4]
          \lgrid{-0.5}{6}{-0.5}{7}
          \dpt{1}{2}\dpt{2}{3}\dpt{3}{3}\dpt{4}{5}\dpt{5}{6}
          \node[font=\scriptsize\bfseries] at (2.5,-1.6){(a)};
        \end{tikzpicture}
        \hspace{0.6cm}
        \begin{tikzpicture}[scale=0.4]
          \lgrid{-0.5}{6}{-0.5}{7}
          \dpt{1}{4}\dpt{2}{2}\dpt{3}{5}\dpt{4}{3}\dpt{5}{4}
          \node[font=\scriptsize\bfseries] at (2.5,-1.6){(b)};
        \end{tikzpicture}
        \end{center}
        (a) is \ans{positive}; \quad (b) is \ans{none}.

  \item \textbf{Match $r$ to the plot.} Assign each $r\in\{\,0.98,\ -0.60,\ 0.05\,\}$ to a plot.
        \begin{center}
        \begin{tikzpicture}[scale=0.4]
          \lgrid{-0.5}{6}{-0.5}{7}
          \dpt{1}{1}\dpt{2}{2}\dpt{3}{3}\dpt{4}{5}\dpt{5}{5}
          \node[font=\scriptsize\bfseries] at (2.5,-1.6){(P)};
        \end{tikzpicture}
        \hspace{0.4cm}
        \begin{tikzpicture}[scale=0.4]
          \lgrid{-0.5}{6}{-0.5}{7}
          \dpt{1}{5}\dpt{2}{3}\dpt{3}{4}\dpt{4}{2}\dpt{5}{3}
          \node[font=\scriptsize\bfseries] at (2.5,-1.6){(Q)};
        \end{tikzpicture}
        \hspace{0.4cm}
        \begin{tikzpicture}[scale=0.4]
          \lgrid{-0.5}{6}{-0.5}{7}
          \dpt{1}{3}\dpt{2}{5}\dpt{3}{2}\dpt{4}{4}\dpt{5}{3}
          \node[font=\scriptsize\bfseries] at (2.5,-1.6){(R)};
        \end{tikzpicture}
        \end{center}
        (P): $r={\color{keyred}\mathbf{0.98}}$; \quad (Q): $r={\color{keyred}\mathbf{-0.60}}$; \quad
        (R): $r={\color{keyred}\mathbf{0.05}}$.

  \item \textbf{Interpret and predict.} A biologist records the number of \emph{days} a plant has grown
        ($x$) and its \emph{height in cm} ($y$). Over days $0$ to $30$, technology gives
        $\hat y=0.8x+5$ with $r=0.97$.
        \begin{enumerate}[label=(\alph*), leftmargin=1.8em, itemsep=4pt]
          \item Interpret the slope $0.8$ in context.
                \ans{each additional day of growth is associated with about $0.8$ cm more height.}
          \item Interpret the $y$-intercept $5$ in context.
                \ans{at day $0$ (the start) the predicted height is $5$ cm.}
          \item Predict the height at $x=20$ days.
                \begin{work}
                  \hat y &= 0.8(20) + 5 \\
                         &= 16 + 5      \\
                         &= 21
                \end{work}
                Height $\approx$ \ans{$21$} cm. Interpolation or extrapolation?
                \ans{interpolation ($20$ is within $0$--$30$)}
          \item Predict at $x=100$ days.
                \begin{work}
                  \hat y &= 0.8(100) + 5 \\
                         &= 80 + 5       \\
                         &= 85
                \end{work}
                Height $\approx$ \ans{$85$} cm. Why is this prediction unreasonable?
                \ans{$100$ days is far outside the data (extrapolation); a real plant stops growing, so it will not keep rising linearly to $85$ cm.}
        \end{enumerate}

  \item \textbf{Correlation vs.\ causation.} A study finds a strong positive correlation between children's
        \emph{shoe size} and their \emph{reading level}. Does a bigger shoe \emph{cause} better reading?
        \ans{no} \; Name a \textbf{lurking variable}.
        \ans{the child's \emph{age} --- older children have bigger feet \emph{and} read better.}

  \item \textbf{Explain.} What two things does the correlation coefficient $r$ tell you about a scatterplot,
        and how does each show up in the value of $r$?
        \ansline{Direction --- the \emph{sign} of $r$ ($+$ rising, $-$ falling); and strength --- the \emph{size} $|r|$.}
        \ansline{$|r|$ near $1$ is a tight linear pattern; $|r|$ near $0$ is little or no linear trend.}
\end{enumerate}
\end{notesbox}

\begin{extensionbox}
\textbf{Which line fits better?} For the five data points below, two candidate lines are proposed:
$y_1=1.5x+0.5$ and $y_2=x+2$. Fill in each line's predictions and its \emph{error} (error $=$ data $-$
prediction) at every point.
\begin{center}
\renewcommand{\arraystretch}{1.3}
\begin{tabular}{|l|c|c|c|c|c|}
\hline
$x$ & $1$ & $2$ & $3$ & $4$ & $5$ \\ \hline
data $y$ & $2$ & $3$ & $5$ & $6$ & $8$ \\ \hline
$y_1=1.5x+0.5$ & $\mathbf{2}$ & $\mathbf{3.5}$ & $\mathbf{5}$ & $\mathbf{6.5}$ & $\mathbf{8}$ \\ \hline
error & $\mathbf{0}$ & $\mathbf{-0.5}$ & $\mathbf{0}$ & $\mathbf{-0.5}$ & $\mathbf{0}$ \\ \hline
$y_2=x+2$ & $\mathbf{3}$ & $\mathbf{4}$ & $\mathbf{5}$ & $\mathbf{6}$ & $\mathbf{7}$ \\ \hline
error & $\mathbf{-1}$ & $\mathbf{-1}$ & $\mathbf{0}$ & $\mathbf{0}$ & $\mathbf{1}$ \\ \hline
\end{tabular}
\end{center}
\begin{enumerate}[label=(\alph*), leftmargin=1.9em, itemsep=6pt]
  \item Total error \emph{size} (add the absolute values of the errors): \; $y_1$: \ans{$1.0$}; \quad
        $y_2$: \ans{$3.0$}.
  \item The line of best fit makes predictions closest to the data overall. Which candidate fits better, and
        how do you know?
        \ans{$y_1$ --- its total error size ($1.0$) is smaller than $y_2$'s ($3.0$), so its predictions sit closer to the data.}
\end{enumerate}
\end{extensionbox}

\begin{spiralbox}
\textbf{Coming next --- the Unit 1 Test, then Unit 2: Quadratic Functions.} You've now met linear functions
from every angle: characteristics, forms, absolute value, piecewise, and fitting a line to data. In Unit~2
the data can \emph{curve} --- a falling-then-rising pattern that no line can capture --- and the curve of
best fit becomes a \textbf{parabola}.
\end{spiralbox}

\end{document}
